% ===========================================================================
% 从通用办公 Agent 到 Research Copilot
% ——WorkBuddy 的能力、腾讯生态与商业化探索
% ---------------------------------------------------------------------------
% 本文件由《WorkBuddy专题研究最终报告_行研版.docx》迁移而来。
% 原模板:HFH-LaTeX-Vorlage (Ilya Zarubin, LPPL-1.3c)
% 已修改:中文编译(XeLaTeX + ctex)、封面中文化、正文替换为报告内容。
%
% 编译方式(重要):
%   Overleaf  : 菜单 -> 编译器 选择 "XeLaTeX",直接编译。
%   本地命令行: latexmk -xelatex main.tex
%   或逐次运行: xelatex main.tex && xelatex main.tex
%
% 注意:切勿使用 pdfLaTeX 编译,否则中文无法显示。
% ===========================================================================
\documentclass[
  12pt,
  a4paper,
  oneside,
  headings=normal,
  numbers=noendperiod,
  listof=totoc,
  bibliography=totoc
]{scrreprt}

% ---------------------------------------------------------------------------
% 中文支持:ctex 包(XeLaTeX 下自动启用 xeCJK,并中文化章节标题/目录等)
% ---------------------------------------------------------------------------
\usepackage{ctex}

% ---------------------------------------------------------------------------
% 元数据
% ---------------------------------------------------------------------------
\newcommand{\MainTitle}{从通用办公 Agent 到 Research Copilot}
\newcommand{\ThesisTitle}{从通用办公 Agent 到 Research Copilot}
\newcommand{\ThesisSubtitle}{基于研究、事实核验、办公编辑与生态闭环的实测}
\newcommand{\AuthorName}{程浩然}
\newcommand{\SubmissionDate}{2026 年 8 月}
\newcommand{\ThesisVersion}{正式版}

\newif\ifIncludeAbstract
\IncludeAbstracttrue
\newif\ifIncludeLegalRegisters
\IncludeLegalRegistersfalse
\newif\ifIncludeAIRegister
\IncludeAIRegistertrue
\newif\ifIncludeAppendix
\IncludeAppendixtrue

% ---------------------------------------------------------------------------
% 语言、字体与页面布局
% ---------------------------------------------------------------------------
\usepackage{microtype}
\usepackage{xcolor}
\usepackage{etoolbox}

\usepackage[
  a4paper,
  portrait,
  left=3cm,
  right=4cm,
  top=2.5cm,
  bottom=2.5cm
]{geometry}

\usepackage{setspace}
\onehalfspacing
\setlength{\parindent}{0pt}
\setlength{\parskip}{6pt}

\setcounter{secnumdepth}{2}
\setcounter{tocdepth}{2}
\setkomafont{disposition}{\normalfont\bfseries\setstretch{1}}
\setkomafont{chapter}{\fontsize{15}{18}\selectfont}
\setkomafont{section}{\fontsize{12.5}{15}\selectfont}
\setkomafont{subsection}{\fontsize{12}{14.5}\selectfont}
\setkomafont{subsubsection}{\fontsize{12}{14.5}\selectfont}
\RedeclareSectionCommand[beforeskip=0pt,afterskip=18pt]{chapter}
\RedeclareSectionCommand[beforeskip=18pt,afterskip=12pt]{section}
\RedeclareSectionCommand[beforeskip=18pt,afterskip=12pt]{subsection}
\RedeclareSectionCommand[beforeskip=18pt,afterskip=12pt]{subsubsection}

% ---------------------------------------------------------------------------
% 页眉页脚、脚注、浮动体与表格
% ---------------------------------------------------------------------------
\usepackage[automark,singlespacing=true]{scrlayer-scrpage}
\clearpairofpagestyles
\ihead{\AuthorName}
\ohead{\MainTitle}
\cfoot{}
\AtBeginDocument{%
  \cfoot{第 \thepage\ 页,共 \pageref*{LastNumberedPage} 页}}
\setkomafont{pageheadfoot}{\normalfont\fontsize{10}{12}\selectfont}
\setlength{\headheight}{15pt}
\pagestyle{scrheadings}
\renewcommand*{\chapterpagestyle}{scrheadings}

\usepackage{chngcntr}
\counterwithout{footnote}{chapter}
\setlength{\footnotesep}{16pt}
\deffootnote[0.5cm]{0.5cm}{0cm}{%
  \raggedright\makebox[0.5cm][l]{\textsuperscript{\thefootnotemark}}}

\usepackage{graphicx}
\usepackage{booktabs}
\usepackage{tabularx}
\usepackage{array}
\usepackage{caption}
\captionsetup{
  format=plain,
  justification=centering,
  singlelinecheck=false,
  font={small},
  labelfont={bf},
  labelsep=colon,
  skip=6pt
}
\counterwithout{figure}{chapter}
\counterwithout{table}{chapter}

\newenvironment{hfhtable}[1][htbp]
  {\begin{table}[#1]\centering\footnotesize\setstretch{1}}
  {\end{table}}
\newcommand{\source}[1]{%
  \par\vspace{3pt}{\footnotesize
    数据来源: #1\par}}
\newcolumntype{Y}{>{\raggedright\arraybackslash}X}

\usepackage[round,authoryear]{natbib}
\setlength{\bibhang}{0.5cm}
\setlength{\bibsep}{6pt}
\AtBeginEnvironment{thebibliography}{\singlespacing}

\usepackage{xurl}
\usepackage[hidelinks]{hyperref}
\usepackage{bookmark}
\hypersetup{
  pdftitle={\MainTitle},
  pdfauthor={\AuthorName},
  pdfsubject={WorkBuddy 专题研究}
}

% ---------------------------------------------------------------------------
% 封面
% ---------------------------------------------------------------------------
\newcommand{\MakeThesisTitle}{%
  \setcounter{page}{1}
  \thispagestyle{empty}
  \begin{center}
    \vspace*{1.2cm}
    {\large 东方证券研究所面试课题 · WorkBuddy 专题研究\par}
    \vspace{0.3cm}
    {\large 专题研究报告\par}

    \vfill
    {\bfseries\fontsize{22}{30}\selectfont\ThesisTitle\par}
    \vspace{0.6cm}
    {\large\ThesisSubtitle\par}
    \vfill

    \begin{tabular}{@{}ll@{}}
      作者: & \AuthorName \\
      日期: & \SubmissionDate
    \end{tabular}
  \end{center}
  \clearpage
}

\makeatletter
\newcommand{\ThesisTableOfContents}{%
  \chapter*{\contentsname}
  \phantomsection
  \addcontentsline{toc}{chapter}{\contentsname}
  \@starttoc{toc}
}
\makeatother

\begin{document}

\MakeThesisTitle

% ---------------------------------------------------------------------------
% 摘要
% ---------------------------------------------------------------------------
\ifIncludeAbstract
  \addchap{摘要}

  本报告要回答一个具体的问题:当研究型知识工作被交给通用办公 Agent,哪些
  环节真正受益,哪些环节反而添乱。我把 WorkBuddy、ChatGPT Work 与 TRAE
  Work 三款产品放进真实工作流里做对照:先是开放式专题研究,再做证据收口式
  的事实核验,最后是直接修改办公文件;随后对 WorkBuddy 的腾讯生态能力单独
  做了一层三层证据链测试,从本地材料处理、真实连接器闭环,一直测到与人工
  路径的摩擦对比。

  结论有三条。第一,WorkBuddy 的能力上限不低:办公文件编辑得分 92.10 居
  三家之首,强模型配置下研究能力 91.25 分逼近 ChatGPT Work 的 92.81 分,
  ima 知识库实现了入库、检索、引用与摘要的全链路。第二,它的表现高度依赖
  底层模型与配置:默认配置下来源质量偏弱,积分消耗波动明显,事实核验环节
  必须有人工复核兜底。第三,生态连接是真实的,但价值按环节分化——知识库
  与文档读写环节确实省时间,PPT 加邮件的分发环节反而比人工更慢,一次实测
  耗时约一小时、干预四次、邮件仍未发出。

  我给出的定位判断是:WorkBuddy 现阶段是\textbf{“带人工复核的研究与办公
  助手”},而不是可以无人值守的办公替身。它能承担材料消化、初稿生成、文件
  编辑与知识沉淀这类前期工作,但口径判断、关键事实核验与对外交付仍需人工
  收口。这个判断贯穿全文,也决定了后面对产品定位、商业化节奏与优化路线的
  讨论。
\fi

\ThesisTableOfContents
\listoffigures
\listoftables

% ---------------------------------------------------------------------------
% 缩写表
% ---------------------------------------------------------------------------
\addchap{缩写表}
\begin{tabular}{@{}p{2.5cm}%
  p{\dimexpr\textwidth-2.5cm-2\tabcolsep\relax}@{}}
  AI   & 人工智能 \\
  LLM  & 大语言模型 \\
  Agent & 智能体:可自主调用工具、完成多步骤任务的程序 \\
  MAX  & WorkBuddy 的高算力增强模式 \\
  HERB & 腾讯混元推出的推理行为评测基准 \\
  MuDABench & 中文企业多文档分析评测集(本项目自建) \\
  ima  & 腾讯 ima.copilot 智能工作台(知识库) \\
  KDocs & 金山文档 \\
  ROI  & 投资回报率 \\
\end{tabular}

% ===========================================================================
% 正文
% ===========================================================================

\chapter{研究背景与核心问题}

\section{研究出发点}

研究型工作不是一个个孤立的问题,而是一条连续的问答链。以写一份行业研究
报告为例:先要判断哪些资料可信、哪些来源值得引用;再逐项核对财务数字,
把结论落进表格和正文;报告交出去之后,新材料、新版本、修改意见会不断
进来,于是又要修订、又要更新口径、又要重新核对一遍。整个过程里,大部分
时间其实花在“找、验、改”上,真正“写”的时间反而是少数。

正因为如此,单轮问答式的模型评测对这类工作几乎没有参考价值。模型能答得
漂亮,不代表它能把一次开放式检索整理成可继续使用的底稿,更不代表它能
记住上一版结论、在追加材料之后正确更新。我考察 WorkBuddy 的方式因此很
直接:把它放进真实的研究流程里,看它在每一个环节是替人省时间,还是给人
添麻烦。

\section{分析框架}

研究按五层递进展开:产品能力、真实工作流价值、可量化价值、付费意愿与
生态放大效应。这五层不是并列关系,而是层层收紧:产品能力是必要不充分
条件,只有通过了真实工作流检验的能力才有商业意义;可量化价值决定了用户
愿不愿意付费;而生态是放大还是缩小这种价值的外部变量。产品与能力的官方
说明见 \citep{workbuddy}。

\begin{figure}[htbp]
  \centering
  \includegraphics[width=0.85\textwidth]{figures/fig1.png}
  \caption{从产品能力到生态协同的分析路径}
  \label{fig:path}
\end{figure}

\section{参照对象的选择}

对照产品选了 ChatGPT Work 与 TRAE Work 两款。ChatGPT Work 是目前把研究
做成品的流程最完整的代表,我用它考察“多步骤研究任务的稳定性”;TRAE
Work 以速度和成本见长,我用它考察“效率维度”。早期考虑过的 Codex 更接近
软件开发场景,Manus 未纳入核心实测,这两项调整都有明确的决策记录,并非
遗漏。选择参照物本身也是研究设计的一部分:参照物决定了我能从差异中读出
什么。

\chapter{研究设计与证据边界}

\section{实验总体设计}

研究由三项竞品测试与一组腾讯生态验证构成,四项实验分别对应工作链的一个
环节:测试一考察“从公开资料形成研究交付”,对应链条的开头;测试二考察
“把结论落回证据、在材料不支持时停止”,对应链条的核验环节;测试三考察
“直接修改办公文件”,对应链条的交付环节。腾讯生态部分单独采用三层证据
链,先做材料处理模拟,再做真实连接器闭环,最后与人工路径比较摩擦。

\begin{figure}[htbp]
  \centering
  \includegraphics[width=0.85\textwidth]{figures/fig2.png}
  \caption{腾讯生态测试的三层证据链}
  \label{fig:layers}
\end{figure}

\section{研究假设}

本研究的每一项实验都对应一个明确的假设。研究不是"跑一遍看看结果",
而是先提出假设、再设计实验去验证或证伪。六个假设如下:

\begin{itemize}
  \item \textbf{H1 工作流假设}:多步骤、跨文件任务的价值高于单轮问答。若
  成立,则 Agent 能完成真实的研究交付;若只有单轮问答强,则 Agent 对研究
  工作的价值有限。验证实验:测试一、测试三与生态测试的 R1/R2。
  \item \textbf{H2 连续性假设}:核心价值在持续更新而非首次生成。研究工作的
  大头是"改"而不是"写",Agent 能否在追加材料后正确更新旧结论,是它能否
  进入日常工作流的关键。验证实验:生态测试 R1 的增量修订、R2 的知识库
  版本区分。
  \item \textbf{H3 用户分层假设}:技术用户从 Agent 获得的增量有限,非技术
  用户因低门槛获得明显价值。验证实验:生态测试的用户分层观察与 MuDABench
  的人工复核需求。
  \item \textbf{H4 生态假设}:生态连接只有在减少步骤、支持读写、支持远程
  触发时才有竞争优势;若只是"接入了很多系统"却增加步骤,则不构成优势。
  验证实验:生态测试 R1/R2(读写)、R3(分发)与企微(远程触发)。
  \item \textbf{H5 商业化假设}:个人用户为便利付费,企业用户为权限、知识库、
  安全与团队协作付费。验证实验:生态测试的 ROI 对比与企业用户观察。
  \item \textbf{H6 专业化假设}:通用 Agent 能完成基础工作,但垂直行业的
  Research Copilot 需要行业模板、长期知识库、来源规范与跟踪机制。验证实验:
  MuDABench 的口径敏感问题。
\end{itemize}

六个假设不是孤立的:它们从"Agent 对单次任务有没有用"(H1)递进到"对
长期工作有没有用"(H2)、"对谁有用"(H3/H5)、"生态放大还是缩小价值"
(H4),最后落在"专业化是不是必经之路"(H6)。后面的每一章结论,都会回到
这六个假设上做一次对照。

\section{执行方式与评分原则}

这一节是整项研究的方法论基础,四条规则保证了对比的公平与可复核。

\textbf{第一,冻结机制。}任务包、提示词、输入材料、输出文件名与保存位置在
正式运行之前冻结,测试过程中不允许改动。这样做的原因很实际:Agent 类产品
迭代快,如果中途调整提示词或更换材料,结果差异就无法归因——你分不清是
平台能力差异,还是任务定义变了。冻结之后,任何改动都必须走版本变更记录,
相当于让三家在同一张卷子上答题。

\textbf{第二,交付物优先。}评分先看交付物本身,再看效率。不能打开的文件、
无法核对的引用、没有写进文件的修改,都不因文字自述得分。这条规则专门
针对一个现象:Agent 很会“说”,把方案描述得头头是道,但交付的文件打不开、
公式是错的。只信交付物,是这类评测最容易守住也最容易被忽视的一条底线。

\textbf{第三,同任务同输入。}三家 Agent 接收同一版本的任务说明;涉及文件
的任务使用同一份输入文件——测试三用同一个 Excel 工作簿,测试二用同一份
唯一事实材料。输入不一致时得出的任何差异都没有意义。

\textbf{第四,运行记录。}每次运行逐项记录时间、积分消耗、失败次数与人工
干预次数。这些记录不直接参与质量评分,却是后面分析成本可预测性、可靠性
与摩擦的原始证据。

\section{证据边界声明}

需要诚实说明本设计的局限:测试三只选了两道题;企业微信受个人账号条件
限制,无法进入真实组织协作场景;三家模型的配置无法完全对齐,ChatGPT
Work 缺少同口径的积分数据。这些限制决定了结论的适用范围——本报告只针对
本次固定任务与运行条件,不把单次结果外推为平台的长期能力。承认边界,
是为了让结论经得起追问。

\chapter{三项测试的总体结果}

\begin{figure}[htbp]
  \centering
  \includegraphics[width=0.85\textwidth]{figures/fig3.png}
  \caption{三项测试的能力矩阵}
  \label{fig:matrix}
\end{figure}

\begin{hfhtable}
  \caption{三平台能力矩阵(主评分)}
  \label{tab:matrix}
  \begin{tabularx}{\textwidth}{@{}Yrrr@{}}
    \toprule
    \textbf{任务} & \textbf{WorkBuddy} & \textbf{ChatGPT Work} & \textbf{TRAE Work} \\
    \midrule
    测试一(研究) & 83.68(HY3) & 92.81 & 88.64 \\
    测试二(核验,平均主评分) & 76.68 & 74.99 & 76.69 \\
    测试三(办公编辑) & 92.10 & 65.05 & 51.80 \\
    \bottomrule
  \end{tabularx}
  \source{三项测试评分结果,口径见第 2 章}
\end{hfhtable}

先说明这张表的读法:三列分数来自不同的评分口径——测试一按五个维度加权,
测试二是两个子任务的平均,测试三是交付物核验。它们之间不能纵向相加,
这也是我没有把三家硬排成一个总分的原因。把口径不同的分数强行合并成一个
名次,看似简洁,实际掩盖了每个平台在不同环节的真实强弱。

在各自的评分口径下,三个横向判断值得注意。其一,ChatGPT Work 在研究任务
中最稳,是唯一在事实、完整性、观点、来源四个维度都没有明显短板的平台。
其二,TRAE Work 的速度优势在测试一中非常突出,总耗时只有 WorkBuddy 的三
分之一左右,但代价是研究完整性的短板。其三,WorkBuddy 的优势集中在办公
编辑——92.10 分大幅领先,这与其在测试一、二中的“中游偏上”位置形成鲜明
对比,说明它的能力不是均匀分布的,评价它必须按场景拆开看。三个测试各自
的原因,在下面三章逐一展开。需要说明的是,这三个测试共同回答了 H1(多
步骤任务的价值):三家都能完成多步骤交付,但完成质量差异显著,而生态测试
的 R3 将证明"多步骤"本身并不自动等于"高价值"——这正是 H1 需要被
精确化的地方。

\chapter{测试一:能力上限不低,但默认配置不稳定}

\section{实验目的}

这一测考察的不是“模型能列举多少事实”,而是“一次开放式检索能否被做成
可继续使用的研究底稿”。合格结果的验收标准有三条:事实可回查、模型的
迭代主线清晰、留下的是数据库、备忘录或汇报材料这类可复用产物,而不是
一段读完就完的长回答。用这个标准的原因很简单:研究者的日常工作是在底稿
上反复加工,而不是每次从零开始生成。

\section{资料、操作与评分}

三家 Agent 接收同一份专题研究提示词,自行检索公开资料并交付五类文件。
WorkBuddy 的主测试使用混元 Hy3,随后用 Auto+MAX 与 DeepSeek V4 Flash+MAX
两种配置复跑,以区分“平台上限”与“模型配置”各自的影响。

评分由四部分构成:任务输出本身、30 条分层事实抽样、12 个核心模型家族的
检查表,以及每家 20 条引用的抽查。证据标准是严格的——一手来源只认官方
公告、官方代码仓库、模型卡、原始论文与技术报告,媒体报道和聚合页可以
作为线索,但不作为事实依据。这个标准直接决定了“来源与引用”维度的得分。

\section{实验结果}

\begin{hfhtable}
  \caption{测试一主测试:维度得分与耗时}
  \label{tab:test1}
  \begin{tabularx}{\textwidth}{@{}Yrrrrr@{}}
    \toprule
    \textbf{平台} & \textbf{总分} & \textbf{事实} & \textbf{完整} & \textbf{观点} & \textbf{来源} \\
    \midrule
    ChatGPT Work & 92.81 & 95.00 & 92.82 & 95.00 & 94.92 \\
    TRAE Work & 88.64 & 93.33 & 81.89 & 82.50 & 95.75 \\
    WorkBuddy(HY3) & 83.68 & 86.67 & 98.33 & 77.50 & 70.25 \\
    \bottomrule
  \end{tabularx}
  \source{测试一评分结果;总耗时:ChatGPT 约 449 秒、TRAE 约 285 秒、WorkBuddy 约 894 秒}
\end{hfhtable}

如表~\ref{tab:test1} 所示,主测试的结论是:ChatGPT Work 以 92.81 分居首,更重要的是四个维度都在
高位,没有明显短板;TRAE Work 88.64 分,速度第一,但研究完整性只有 81.89,
比第一名低了近 11 分;WorkBuddy(Hy3)83.68 分,研究完整性反而是三家最高
的 98.33,说明它愿意也擅长铺开较宽的材料面,但来源与引用只有 70.25——
抽样检查中一手来源比例偏低,较多依赖媒体报道、聚合页与第三方模型卡。
这意味着它产出的参数、发布日期、定价与排名等信息,都需要更高强度的人工
复核才能采用。

\begin{figure}[htbp]
  \centering
  \includegraphics[width=0.85\textwidth]{figures/fig4.png}
  \caption{测试一主测试结果}
  \label{fig:test1-main}
\end{figure}

\section{换模型复测:平台与模型的分离}

主测试只能说明“默认配置下的 WorkBuddy”。为了分离平台与模型的影响,
我把同样的任务换到 Auto+MAX 与 DeepSeek V4 Flash+MAX 两种配置下重跑,
结果改变了对它的判断:Auto+MAX 得 91.25 分,逼近 ChatGPT Work;DeepSeek
V4 Flash+MAX 得 90.00 分,耗时约 290 秒,接近 TRAE Work 的速度水平。
上下文消耗同样值得注意:DeepSeek 配置只用到 142.2k/1000k,余量明显大于
Auto 配置的 150.5k/168.0k,说明它在长任务中更“省”。

\begin{figure}[htbp]
  \centering
  \includegraphics[width=0.85\textwidth]{figures/fig5.png}
  \caption{WorkBuddy 换模型复测结果}
  \label{fig:test1-rerun}
\end{figure}

\section{本章判断}

我的判断是:WorkBuddy 的平台上限不低,但用户实际获得的质量,取决于底层
模型、MAX 模式、上下文容量与工具链的组合。它不是一个“固定低分”的产品,
而是一个“能力高度依赖配置”的产品。对采购方来说,这意味着只看产品入口
是不够的——必须问清楚默认路由到哪个模型、上下文多大、MAX 模式怎么开、
积分怎么消耗,否则同一平台会表现出完全不同的质量与成本结构。

\chapter{测试二:会引用证据,不等于结论正确}

\section{实验目的}

事实核验考察的不只是“搜到链接”,而是三件更难的事:把一句声明拆解成可
核验的原子事实、识别证据的边界在哪里、确认证据的位置是否真实存在。本
测试还顺带观察一个问题:更长的回答、更多的积分消耗,是否真的带来更可靠
的结论。

\section{资料与操作}

测试二由两个子任务构成。第一组采用 Factcheck-GPT 与 FactBench 共 12 题,
覆盖可核验性判断、声明事实性与文档级核验 \citep{factcheckgpt};第二组
采用 Google DeepMind 的 FACTS Grounding 公开样本 9 题,按短、中、长上下文
各取三题 \citep{factsgrounding}。运行期间禁止联网,给定的唯一事实材料是
答题依据。两个子任务考察的能力并不相同:前者暴露声明拆解与事实性判断,
后者暴露上下文边界控制、证据定位与复杂指令遵循。

\section{Factcheck-GPT 子任务}

\begin{hfhtable}
  \caption{测试二:两个子任务的主评分与资源消耗}
  \label{tab:test2}
  \begin{tabularx}{\textwidth}{@{}Yrrrrr@{}}
    \toprule
    \textbf{平台} & \textbf{Factcheck 主评分} & \textbf{FACTS 主评分} & \textbf{平均主评分} & \textbf{总耗时(秒)} & \textbf{总积分} \\
    \midrule
    WorkBuddy & 60.02 & 93.34 & 76.68 & 3622.80 & 498.70 \\
    ChatGPT Work & 59.55 & 90.43 & 74.99 & 2499.75 & 未记录 \\
    TRAE Work & 57.50 & 95.88 & 76.69 & 3292.65 & 413.31 \\
    \bottomrule
  \end{tabularx}
  \source{测试二评分结果;ChatGPT Work 未记录积分消耗}
\end{hfhtable}

如表~\ref{tab:test2} 所示,两个子任务放在一起看,没有一家在两个任务上
同时领先:Factcheck 任务中 WorkBuddy 综合得分最高(60.02),FACTS 任务中
TRAE Work 领先(95.88),ChatGPT Work 则两项都居中。得分差异背后是能力
结构的不同——WorkBuddy 强在文档级核验与声明事实性,ChatGPT Work 强在
可核验性判断,TRAE Work 强在材料边界控制。此外,文档级任务预期拆出 60
条原子事实,ChatGPT 输出 27 条,WorkBuddy 输出 29 条,TRAE 输出 36 条:
拆分数量最多并不自动等于质量最高,关键在于是否贴合标准事实单元,但 TRAE
更积极的拆分确实带来了更高的文档级核验得分。

\section{FACTS Grounding 子任务}

这一组 TRAE Work 表现最强,主评分 95.88,短中上下文任务的材料边界控制
很好,总耗时也最低;风险在于长上下文任务中出现了不存在的段落编号,这会
直接降低人工复核的效率。WorkBuddy 93.34 分,长上下文覆盖能力较强,但
同样存在证据映射问题,比如部分段落编号错误、CSV 行错位。ChatGPT Work
90.43 分,证据引用最为保守,倾向于使用真实存在的段落编号,降低了虚构
证据位置的风险;但它在 FACTS\_004 中误读了题目里的关键约束,在另一题中
违反了“不得使用外部工具”的指令——它的短板不在输出结构,而在复杂条件下
的约束识别。

\begin{figure}[htbp]
  \centering
  \includegraphics[width=0.85\textwidth]{figures/fig6.png}
  \caption{两类事实核验任务的主评分}
  \label{fig:test2}
\end{figure}

\section{积分与质量:一个被忽视的量化发现}

把两个子任务的积分记录放在一起,得到一个值得单独讨论的发现:WorkBuddy
的积分消耗波动明显大于 TRAE——单题积分从 4.47 到 54.97,极差超过 12 倍,
变异系数约 67\%;而积分与质量的相关性接近零(Factcheck 子任务中 WorkBuddy
为 -0.083,TRAE 为 -0.134)。也就是说,多花积分并不带来更高的核验质量,
积分波动更可能反映的是任务路径、检索过程与平台内部资源调度的差异。这个
发现的意义在于:成本不可预测是 WorkBuddy 的真实短板,而且它不能靠“多花
钱买质量”来弥补。

\section{本章判断}

我的判断是:事实核验类任务中不存在绝对领先的平台,三家各有所长,也各有
失分点。WorkBuddy 质量均值略高、长材料覆盖较强,但代价是积分消耗高且
波动大;它可以承担声明拆解、证据整理与不确定性暴露,但关键数字、行业
口径与对外结论必须由人确认——证据可复核,不等于结论一定正确。

\chapter{测试三:办公文件是否真正改好}

\section{实验目的}

前两项测试都以 Markdown 交付,容易让“会写说明”掩盖“未真正改文件”。办公
编辑测试直接把考察对象换成文件本身:Agent 能否在理解现有工作簿、不破坏
其他内容的前提下修复公式或创建图表,并交回一份可继续使用的 Excel?评分
不再看描述,而是直接检查输出文件能否打开、公式是否正确、图表对象是否
真实存在。

\section{资料与操作}

资料取自 SpreadsheetBench 2 \citep{spreadsheetbench},受时间限制选取两道
题:一道是财务模型的硬编码与公式修复,一道是两月数据合并并生成折线图。
三家使用同一份输入文件与固定的输出名,交回后可被统一核验。

\section{实验结果}

第一题,WorkBuddy 与参考答案的一致率 99.53\%、公式一致率 97.08\%,ChatGPT
Work 为 99.67\% 与 97.77\%——两家几乎打平,WorkBuddy 略低但高度接近;
TRAE Work 的输出文件在这一题出现 XML 解析错误,按评分规则视为无法可靠
复核,直接触发交付门槛。

第二题的差距被彻底拉开。WorkBuddy 与 TRAE 都检测到 1 个真实折线图、4 个
系列、缺失值显示为空缺;而 ChatGPT Work 的变更说明声称已经创建图表,文件
内部却没有真实的图表对象,数据区还残留 150 个无效的等号公式。同样是
“说做了”,只有 WorkBuddy 和 TRAE 真正交出了可打开、可核验的图表文件。
效率与成本方面,TRAE 最快(约 446 秒),WorkBuddy 最慢(约 1473 秒)且两题
消耗约 315 积分;WorkBuddy 第二题的上下文用到 117.3K/168K,说明复杂办公
任务会快速消耗上下文与额度。

\section{本章判断}

我的判断是:办公 Agent 的分水岭不在“会不会分析任务”,而在“能否稳定交付
真实、可打开、可继续编辑的文件”。WorkBuddy 在这一项上综合最稳,适合重
交付、可复核的任务;ChatGPT Work 的公式审计很强,但图表对象落盘不可靠;
TRAE Work 极快且成本低,但文件结构损坏风险让它存在高风险失效点。这也是
WorkBuddy 最接近日常办公价值的部分——前提是它愿意把速度与稳定性补上。
\chapter{腾讯生态测试:三层证据链的设计与验证}

\section{实验目的}

前三项测试说明的是单任务能力,不能证明“腾讯生态”带来的增量。生态价值
应该体现在软件之间:资料进入文档云与知识库之后,能否被继续读取、修改与
引用;研究成果能否再进入 PPT 与邮件。只有跑通整条链,才能把模型能力、
连接器能力与真实工作流的摩擦区分开来——这正是本模块与前三项测试的本质
区别。

\section{三层证据链的设计逻辑}

\subsection{L1:材料能力模拟}

第一层在本地完成,材料是腾讯场景的典型形态(项目文档、会议记录、沟通
消息、PR 状态),要求生成初版报告并做增量修订。这一层刻意排除网络与
连接器干扰,只回答一个问题:WorkBuddy 能不能消化这类材料、能不能在追加
材料后正确更新结论。它证明材料处理能力,但不证明生态价值。

\subsection{L2:真实工具闭环}

第二层接入真实连接器,做三条链路:R1 文档云创建、读回、修订、评论与
分享;R2 ima 知识库的入库、检索、引用、摘要与版本修订;R3 从报告生成
五页 PPT 并尝试通过邮箱发送。这一层回答的问题是:生态连接能不能真实地
读写、沉淀与分发。

\subsection{L3:摩擦对比}

第三层把 R3 的真实执行与人工路径做摩擦对比,比较耗时、人工介入次数与
交付结果。这一层回答的是最实际的问题:就算链路能跑通,它比人自己动手
是更快还是更慢?值不值得用?

三层依次收紧:L1 排除干扰、L2 引入真实环境、L3 引入人工基线。任何一层
单独存在都会得出偏颇的结论——只跑 L1 会高估生态,只跑 L3 又无法定位
问题出在模型还是连接器。

\chapter{L1:材料理解能力成立,口径仍需人工收口}

\section{HERB:跨材料分析与增量修订}

HERB 材料包模拟了一个真实项目的多来源状态:项目文档、会议记录、沟通
消息与 PR 状态交织在一起,信息互相印证也互相冲突。任务要求生成初版报告,
并在追加材料后完成增量修订。四组配置跑下来,DeepSeek V4 Flash 的初版与
修订版均得满分;Hy3 得 96 分,开启 MAX 后升至 98.5 分——它能识别产品
方向变化、PR 状态与材料冲突,并完成修订;初版的小幅失分来自一处敏感性
分析的漏检,而 MAX 模式对此有明显补救作用。

\begin{figure}[htbp]
  \centering
  \includegraphics[width=0.85\textwidth]{figures/fig8.png}
  \caption{HERB 四种模型组合的两阶段均分}
  \label{fig:herb}
\end{figure}

HERB 还有一个细节值得单独说:面对一个无法回答的问题,DeepSeek 正确选择了
拒答,而不是编造一个看似合理的答案。“知道什么时候该停下”是研究型任务的
基本素养,也是很多模型评测不会专门考、但真实工作里天天遇到的能力。

\section{MuDABench:口径敏感问题的实证}

MuDABench 取六道中文企业多文档题,最终只跑了 DeepSeek 一组。平均 71.50
分,两道满分(资本充足率前三、第一大股东前三),四道出现口径分歧;但六道
题的证据可复核性全部满分——文件、页码、原文引用都极其扎实。失分不是
来自事实错误或证据缺失,而是来自“结论与锚点口径不一致”。

四处分歧逐一说明:营业总成本一题,锚点把银行也计入样本,但银行利润表
并没有“营业总成本”科目,只有营业支出,模型因此只统计了非银公司,并明确
标注了这层风险;资本增长率一题,锚点用的是资本充足率的变化率(量级在
0.01\%),模型用的是更符合市场习惯的资本净额增速(7\%—10\%);异常公司判定
一题,模型的结论对均值的定义方式高度敏感,剔除异常值前后结论完全不同;
ESG 一题,两份材料里没有股票代码,模型按照“不得外部补全”的规则拒绝了
这两家。这些分歧有一个共同点:模型的逻辑在给定材料内自洽,出错的是
“该用什么口径解读材料”这件事,而这恰恰是人需要介入的地方。

\begin{figure}[h]
  \centering
  \includegraphics[width=0.85\textwidth]{figures/fig9.png}
  \caption{MuDABench 六题得分}
  \label{fig:mudabench}
\end{figure}

我的判断是:模型并没有乱编,甚至主动暴露了不确定项,但“采用哪种业务口径”
最终必须由人来决定。这正是“带人工复核闭环的研究助手”定位在能力层的最
直接证据——它适合做研究,不适合无人值守的深度研究。

这个判断可以从三个层面展开。第一,失分的性质决定了问题的层级。六道题里
证据可复核性全部满分,意味着模型的检索、定位与引用能力没有问题;失分
全部集中在"结论与锚点口径不一致",意味着问题不在"找不到事实",而在
"如何解读事实"。前者是执行层能力,后者是判断层能力,两层是分开的。

第二,模型在口径分歧中的行为是"可复核的坚持"。它在营业总成本一题里
明确标注了"银行利润表没有该科目"的风险,在资本增长率一题里选择了更符合
市场习惯的算法,在 ESG 一题里因材料缺失而拒绝作答——每一步都有明确的
理由,不是随机出错。这恰恰是它适合与人工配合的原因:它会把自己的不确定
暴露出来,人工只需要在这些暴露点上做裁决,而不是从头到尾复查。

第三,这组结果直接支持 H6 专业化假设。通用模型之所以在口径敏感题上失分,
不是因为笨,而是因为它缺少"行业默认用什么口径"这类背景知识。一个面向
金融行业的 Research Copilot,应该预置"营业总成本 vs 营业支出""资本净额
增速 vs 资本充足率变化"这类行业规则,让模型在解读材料时就有正确的默认
值,而不是事后靠人工纠偏。这也是从"通用 Agent"走向"垂直行业 Agent"
最具体的理由:MuDABench 的四道分歧题,就是行业模板必要性的四张成绩单。

\chapter{L2:知识库闭环是最扎实的正面结果}

\section{R1:文档写回与平台偏差}

R1 完成了文档云创建、写回、读回、增量修改、评论与分享的全链路,判据
8 项全部达成,但过程约 74 分钟,效率明显偏低。更值得注意的是一个平台
层面的发现:这次实测实际调用的是金山文档(KDocs),而不是腾讯文档——
连接器页面显示“已连接”,并不保证对应的工具在会话中必然出现。因此 R1
只能作为“文档云闭环”的证据,不能表述为“腾讯文档闭环”。这条观察本身
就是研究的一部分:连接器的信任度,与工具实际暴露的确定性,是两个不同的
问题。

\section{R2:ima 知识库全链路}

R2 使用 ima 知识库完成了四步信息流:16 份 PDF 入库、网页导入、初版与
修订版的增量管理、以及检索、引用与摘要。除建库这一步需要手动完成外,
链路基本走通,全程约 20 分钟,仅一次人工介入。对经常阅读大量报告的研究
人员来说,这项能力确实能省下反复翻文件的时间——16 份年报如果靠人逐份
翻阅并定位数据,通常需要数小时。

\begin{figure}[htbp]
  \centering
  \includegraphics[width=0.85\textwidth]{figures/fig10.png}
  \caption{ima 知识库的解析覆盖情况}
  \label{fig:ima}
\end{figure}

\section{一个必须正视的静默错误}

但 R2 也暴露了一个隐蔽的风险。步骤四检索“资本充足率前三”时,模型给出
建行、上海农商行、交行的排序,与 L1 锚点(工行、建行、农行)不一致。
排查后确认根因:工行与农行的年报已经入库,但当时异步解析尚未完成,这两
份材料根本没有进入对比范围。模型在回答中如实标注了“只比较了已解析文件”,
做法是负责的,但结论在事实上仍然错了。

这个案例的价值在于它指出了知识库工具的一个结构性缺陷:**入库成功、可检索
与结论完整是三件不同的事**。入库成功只代表文件进了存储,可检索只代表
解析已完成,而结论的完整性取决于全部相关材料是否都进入了推理。如果工具
不展示解析覆盖率,用户就无法察觉结论是基于不完整样本得出的——这就是
“静默缺样”。它再次印证了人工复核点的必要性:连“最顺滑可用”的知识库
环节也逃不掉。

\section{本章小结}

L2 证明了生态的“数据层”与“执行层”是真实可用的:文档可以读写,知识库
可以沉淀、检索并支撑创作。但可用性高度不均——知识库环节顺滑,文档环节
低效,分发环节(见下一章)直接失败。生态连接是真实的,价值却必须按环节
评估。

把三个任务的结论放在一起,可以提炼出四点。第一,链路完整性存在明显梯度:
R2 是唯一"零摩擦"走完的闭环,从入库到检索到创作一气呵成;R1 链路虽通
但耗时 74 分钟,效率上不具备日常可用性;R3 则直接失败。梯度不是随机分布,
而是与"工具成熟度"正相关——知识库产品打磨最久,所以最顺;分发链路
依赖最重的附件通道,所以最脆。

第二,R1 的平台偏差(kdocs 而非腾讯文档)揭示了一个被忽略的可靠性维度:
"连接器显示已连接"不等于"工具在会话中必然可用"。对企业采购而言,这
意味着必须用真实任务验证连接深度,而不能信任配置页面的状态;对评测而言,
这意味着每个生态结论都要记录实际后端,否则会把平台 A 的能力记到平台 B
头上。

第三,静默错误是比显式失败更值得警惕的风险。R3 失败会立即暴露,用户
马上知道要人工介入;R2 的错误则藏在"看似正常的结果"里——模型给出了
排名,甚至标注了材料范围,但结论仍然是错的。知识库工具如果不提供解析
覆盖率,这类错误就无法被用户察觉,这是生态产品在"可信度"上最需要补的
一块。

第四,这组结果同时回应了 H2 与 H4 两个假设。H2 连续性假设得到验证:R1
的增量修订与 R2 的版本区分都成功完成,证明 Agent 能承载"持续维护"型
工作,这正是研究工作的核心形态。H4 生态假设则被精确化:生态连接在"减少
步骤"的环节(R2 检索替代翻文件)成立,在"增加步骤"的环节(R1 低效、
R3 失败)不成立——生态价值不是整体的,而是分环节的,这为下一章的摩擦
对比埋下了伏笔。

\chapter{L3:PPT 与邮件分发环节的失败案例}

\section{R3:一次完整的失败过程}

R3 的意图是测通完整办公链路:依据报告生成五页 PPT,再通过 QQ 邮箱发送
出去。结果这是一次教科书式的失败过程,分三个阶段。第一阶段,PPT 生成
的官方技能在本地运行时依赖装不上(安装过程报错、依赖解析失败、原生依赖
缺失),折腾约十分钟后放弃,换用替代方案;第二阶段,PPT 最终生成成功,
但内容被反复压缩,呈现质量完全不合格;第三阶段,邮件附件需要转成 base64,
文件过大导致工具调用多次无返回,反复压缩(从 1.6MB 压到 18KB)仍然发不
出去。

任务总共持续约一小时,人工干预四次,最终邮件未发出,零交付。手工制作
PPT 并发送邮件,估计只需要 10—20 分钟(估算值,非同条件实测)。差距超过
三倍,而且自动化路径没有任何产出,结论已经不需要更多解释。

\section{根因归属:问题在工具链,不在模型}

这次失败最容易的归因是“模型不行”,但细看记录会发现不是:附件参数限制、
技能依赖安装失败、工具调用超时——这些问题换更强的模型依然存在。更严重
的是 Agent 缺乏止损机制:依赖重试了四次、附件重试了三轮、工具调用无返回
时仍沿着同一路径盲目重试,直到用户干预。也就是说,根因在平台工程层面,
而不是模型能力层面。这个区分很重要:如果归因错误,优化方向就会跟着错。

\section{本章判断}

\begin{hfhtable}
  \caption{分发环节的摩擦对比(R3 实测 vs 人工估算)}
  \label{tab:l3}
  \begin{tabularx}{\textwidth}{@{}YYY@{}}
    \toprule
    \textbf{指标} & \textbf{WorkBuddy(实测)} & \textbf{人工(估算)} \\
    \midrule
    总耗时 & 约 60 分钟,未完成 & 10—20 分钟,完成 \\
    人工介入 & 4 次干预 + 1 次确认 & 0 \\
    完成状态 & 失败,0 交付 & 成功 \\
    失败与重试 & 附件重试 3 轮全部失败 & 通常没有 \\
    \bottomrule
  \end{tabularx}
  \source{人工数据为估算区间,仅用于量级对比}
\end{hfhtable}

如表~\ref{tab:l3} 所示,我的判断是:生态覆盖广与工作流好用是两回事,这次
失败把这条界限划得非常清楚。从表面看,问题出在附件传输和技能依赖这些
技术细节上,但它们共同指向一个更深层的结论:连接器的存在,不等于连接器
可用;可用的定义,不只是"能调用",而是"能在合理时间内稳定完成一件事"。

这个失败也修正了我对生态价值的理解。生态价值不是由连接器数量决定的,
而是由"它替用户省下了多少真实动作"决定的。R2 知识库之所以成立,是因为
它把"翻开 16 份年报逐份找数据"压缩成了几次检索;R3 之所以失败,是因为
它把"做 PPT 加发邮件"这件本来 10—20 分钟能做完的事,拖成了一个小时还
没有结果。同样是生态连接,一个减摩擦,一个增摩擦,价值方向完全相反。
因此,评估任何生态功能,都应该先问一句:它让用户少做了什么?而不是它
接入了多少个系统。

对用户而言,现阶段让 Agent 独立完成“做 PPT 并发邮件”的任务,远不如
自己动手。但这不意味着生态路线本身有问题,而是说明分发环节的产品化程度
还远未成熟——它需要工程层面先把附件通道、依赖预置和止损机制做好,才有
资格谈"替代人工"。

\chapter{成本与使用方式:强模型应该配置在哪里}

测试一告诉我们强模型能明显提升研究质量,测试二告诉我们更多积分不一定
提高核验得分。这两组证据看似矛盾,其实指向同一个事实:WorkBuddy 的产出
质量与资源消耗之间没有稳定的对应关系。质量取决于配置,成本取决于路径,
两者是两条独立的曲线。把这两条曲线放在一起看,一个合理的做法浮出水面:
不是全量开启 MAX,而是按风险分级——把强模型花在最需要判断力的环节,
把便宜配置留给重复性工作。

成本结构本身也值得展开。测试二里,WorkBuddy 的积分变异系数约 67\%,单题
积分从 4.47 到 54.97,极差超过 12 倍;TRAE Work 的变异系数只有约 42\%。
也就是说,在 WorkBuddy 上完成同一难度任务,费用可以差出一个数量级,而
用户在下单前并不知道会落在哪一端。更麻烦的是,积分与质量的相关性接近零
(Factcheck 子任务中 WorkBuddy 为 -0.083,TRAE 为 -0.134),多花的积分
并不带来更高的核验得分。对个人用户来说,这意味着"贵"不等于"好";
对企业来说,这意味着预算无法被预测,也就无法被管理。

上下文消耗是另一个被低估的成本维度。测试一中,DeepSeek 配置只用到
142.2k/1000k,余量明显大于 Auto 配置的 150.5k/168.0k;测试三第二题更是
冲到 117.3K/168K。复杂办公任务会快速消耗上下文,而上下文不足往往意味着
任务中途降级或失败——这正是梯度路由需要提前感知的资源信号。

\begin{hfhtable}
  \caption{模型路由建议}
  \label{tab:routing}
  \begin{tabularx}{\textwidth}{@{}Yrr@{}}
    \toprule
    \textbf{任务阶段} & \textbf{建议配置} & \textbf{原因} \\
    \midrule
    整理与分类 & 混元 Hy3 或基础模型 & 任务重复,错误易被后续发现 \\
    跨材料分析 & DeepSeek 或强模型 & 需处理冲突与长上下文 \\
    正式交付前核验 & 强模型 + 人工复核 & 关键事实不能只靠模型自检 \\
    \bottomrule
  \end{tabularx}
  \source{基于测试一、二结果的配置建议}
\end{hfhtable}

但梯度控制有一个前提:平台必须能提前告知大致成本,并在上下文不足、连接
器失败时给出明确反馈。如果做不到这一点,梯度控制就只是用户个人的经验
积累,而不是稳定的产品能力。测试二里 67\% 的积分变异系数,正是这个前提
缺失的量化证据——用户连“这单大概花多少”都猜不准,自然不敢把重要任务
交给它。更进一步说,成本可预测性本身就是产品的一部分:它决定用户敢不敢
把高频任务交给 Agent,也就决定了 Agent 能不能从"偶尔一用的玩具"变成
"每天依赖的工具"。

\chapter{用户分层与付费判断}

\section{五类用户的实测画像}

个人用户最容易感知的价值是少翻文件、少做重复整理,ima 知识库正属于
此类;而 PPT 加邮件的失败体验,则会直接抵消付费意愿。对熟悉编程的研究员
来说,WorkBuddy 的增量价值有限——很多操作可以用脚本替代;对不会编程的
知识工作者,它确实降低了使用门槛,但要求用户具备基本的核验能力,因为
模型会给出看似确定、实际有误的答案。中小团队更看重共享知识库与用量管理;
大型企业则关注权限、安全、审计与企业微信入口——而这三项,恰恰是本项目
未能充分验证的部分。

\begin{hfhtable}
  \caption{不同用户的价值判断}
  \label{tab:users}
  \begin{tabularx}{\textwidth}{@{}YYr@{}}
    \toprule
    \textbf{用户类型} & \textbf{核心价值} & \textbf{付费判断} \\
    \midrule
    个人知识工作者 & 少翻文件、重复整理提效(知识库) & 标准版 ROI 为正;分发场景为负 \\
    技术型研究员 & 增量价值有限,可脚本替代 & 低,倾向免费版或工具组合 \\
    非技术型研究员 & 低门槛 + 需复核闭环 & 中高,标准/高级版 \\
    中小团队 & 共享知识库 + 用量管理 & 高,企业版(¥198/人/月) \\
    大型企业 & 权限、安全、审计、企微入口 & 价值主张最强,稳定性门槛高 \\
    \bottomrule
  \end{tabularx}
  \source{基于本报告全部实测的用户价值判断}
\end{hfhtable}

\section{分场景的 ROI 测算}

按研究员时薪约 150 元估算(以月薪 2.5 万、月工时 170 小时折算),知识库
检索类任务单次净价值约 +240 至 530 元(20 分钟 vs 人工 2—4 小时),报告
生成与修订类约 +300 至 740 元,会议纪要类约 +80 至 240 元;而 PPT 加分发
场景是明确的负价值——约 -150 至 -220 元(60 分钟失败 vs 人工 10—20 分钟
完成)。也就是说,一个每月做十次以上检索类任务的个人用户,标准版(连续
包月 70 元)的 ROI 显著为正;而主要依赖分发场景的用户,基本不会付费。

企业侧的账更清楚:SaaS 版约 198 元/人/月,一个十人团队月成本约 1980 元,
需要每个成员每月稳定节省两小时以上的工作时间才能覆盖——在当前的稳定性
水平下,这个条件很难满足。价值主张最强的其实是大型企业(权限、安全、
审计、生态整合),但恰恰是它们对稳定性要求最高,而稳定性正是当前短板。

\section{本章判断}

综合上述测算,我对 WorkBuddy 的用户分层与付费前景做出以下判断。

第一,个人市场的付费逻辑是“省时间换订阅”,且只在特定环节成立。个人用户
愿意为知识库检索这类任务付费,是因为单次净价值为正、且发生频率高——每月
十次以上检索就能覆盖标准版月费。但同一个用户的付费意愿可以被一次失败的
分发体验直接摧毁:当用户意识到"交给它反而更慢"时,前面积累的信任会
整体归零。因此个人市场的关键是保证正价值场景的高可用,而不是追求场景
覆盖的广度。

第二,企业市场的机会与门槛同样清晰。中小团队的核心诉求是共享知识库与
用量管理,这些是 WorkBuddy 已经具备的能力,且定价(约 198 元/人/月)在
可接受区间;大型企业的价值主张最强——权限、安全、审计与企业微信入口
都是中小企业给不了的差异化,但恰恰是这三项本项目未能验证,也是稳定性
要求最高的部分。企业市场能不能起量,取决于平台工程修复的进度,而不是
销售力度。

第三,也是最重要的一点:现阶段把 WorkBuddy 定义为"统一办公入口"为时
过早。入口层的企业微信个人用户根本够不到,数据层与执行层虽通,分发层
又不可用,三条链路里只有知识库一条是完整闭环。一个还没有跑通一半链路
的产品,贸然以"入口"自居,只会放大用户对缺失环节的期待,然后以失望
收场。更稳妥的商业化路径,是从一个能打的场景(知识库与检索)出发,把
单点体验打磨到无感,再逐步扩展到文档修订、纪要整理,最后在稳定性达标
后进入企业市场。这个顺序与"草房子先装修再出租"是同一个逻辑:先证明
能住人,再谈当办公楼。

\chapter{从 WorkBuddy 到 Research Copilot:优化路线}

做完实验,我不再把 Research Copilot 理解为“更聪明的聊天窗口”,而是一套
完整的工作流:知道哪些文件没有解析、何时应该停下来问人、记住研究口径与
来源范围、维护观点的版本。这条路线不是凭空设计的,每一步都对应一个实测
缺陷。

\begin{figure}[htbp]
  \centering
  \includegraphics[width=0.85\textwidth]{figures/fig12.png}
  \caption{WorkBuddy 的后续优化顺序}
  \label{fig:roadmap}
\end{figure}

\textbf{第一步(短期)解决失控问题。}失败止损机制,对应 R3 里依赖重试四次、
附件重试三轮却仍沿原路径走下去的问题;参数体量预检,对应附件过大导致
工具调用无返回的问题;常用技能依赖预置,对应 PPT 生成技能装不上的问题。
这三个根因,正好覆盖了本轮失败的主要直接原因。

\textbf{第二步(中期)提升结果可信度。}展示知识库解析覆盖率,对应 R2 的
静默缺样——用户需要知道结论基于多少材料;优先发送链接而非硬塞附件,
绕开 base64 附件通道的硬阻塞;把人工确认集中在高风险节点,而不是让用户
全程盯守。这一步的目标是把“需要复核”从被动救火变成主动设计。

\textbf{第三步(长期)才是行业化。}固定字段、来源白名单、口径模板与观点
版本管理,对应 MuDABench 的四题口径分歧——通用模型之所以在这些题上失分,
是因为缺少行业化的判定规则。到这一步,产品形态就接近真正的垂直行业
Research Copilot 了。

支撑这条路线的是一个护城河判断:模型可以被替换(本项目全程使用外部模型,
平台本身没有模型壁垒),连接器目前又不够稳定,那么真正能沉淀下来的,只有
知识库数据、行业模板与长期工作记录这三样东西。它们都依赖长期使用积累,
因此时间窗口既是 WorkBuddy 的护城河建造期,也是被竞品追赶的窗口期。

如果用一句话总结现状,就是之前说的“草房子”:能住,说明方向正确;漏风、
工具不好用、线路不稳,也都是真实存在的问题。当前最重要的不是继续增加
房间,而是把常用链路修好——装修工程量,直接决定商业化能走到哪一步。

\chapter{结论}

\section{已成立的优势}

四项优势都有实测数据支撑,它们共同回答了"WorkBuddy 为什么值得关注"。

一是材料消化与证据定位能力强。HERB 两阶段满分,初版报告与追加材料后的
修订全部正确,还能在材料不足时正确拒答;MuDABench 的证据可复核性六题
全满分,文件、页码、原文引用都经得起回查。这意味着在"把多来源材料整理
成可追溯的结论"这个研究最基础的动作上,它的能力是稳定可靠的。

二是办公文件编辑能力领先。测试三 92.10 分,大幅领先第二名(65.05 分),
能真实修复 36 处公式问题并生成可打开的图表对象,而不是只交一份修改说明。
这是三家中唯一在"交付真实、可打开、可继续编辑的文件"维度上表现稳定的
平台,也是它最接近日常办公价值的部分。

三是知识库闭环可用。ima 实现入库、检索、引用与摘要的全链路,16 份 PDF
的处理替代了人工数小时的翻阅——这是腾讯生态里唯一被完整验证的正价值
环节,也是"生态能沉淀知识"的直接证据。

四是强模型配置下研究能力接近第一梯队。Auto+MAX 复测 91.25 分,逼近
ChatGPT Work 的 92.81 分;DeepSeek V4 Flash+MAX 得 90.00 分且速度接近
TRAE Work。这说明平台的工程与工具链本身不是瓶颈,只要配置到位,它的
上限不输给任何对手。

\section{明确的不足}

五项不足同样有据可查,它们共同决定了 WorkBuddy 目前只能做"助手"而
不能做"替身"。

一是模型配置敏感。Hy3 基线下来源与引用维度只有 70.25 分,换到 Auto+MAX
后整体升到 91.25 分——同一平台、同一任务,质量可以相差近 8 分。用户
的实际体验完全取决于默认路由到了哪个模型,而这种路由对用户是黑箱。

二是成本不可预测。积分变异系数约 67\%,单题积分从 4.47 到 54.97,且积分
与质量相关性接近零。用户既无法预估单次账单,也无法靠多花钱买到更可靠的
结果,成本成为了一笔无法管理的支出。

三是连接器与技能稳定性不足。通道崩溃、依赖装不上、附件传输硬阻塞,均属
平台工程问题,与模型无关——换更强的模型,这些问题一个都不会消失。生态
覆盖面广是优势,但稳定性不均正在抵消这个优势。

四是 Agent 缺乏止损与反思机制。一次失败的任务需要用户四次干预,Agent
会沿着同一条失败路径反复重试,直到用户强制打断。这意味着"全程盯守"
是当前使用 WorkBuddy 的隐性成本,也直接否定了"无人值守"的可能。

五是知识库存在静默缺样风险。入库不等于解析完整,工行与农行的年报缺席
直接导致了一次错误的排名结论,而用户如果不逐份核对根本无从察觉。这个
风险比显式的错误更危险,因为它不会报错。

\section{未来的整改方向}

整改方向不是一份愿望清单,而是把每个实测缺陷映射到对应工程动作之后得到
的优先级顺序。

短期优先修复三个根因,直接消除本轮失败的主要来源:一是建立失败止损机制,
让 Agent 在连续重试无果时改变策略或停下来询问,而不是反复撞同一堵墙;
二是增加参数体量预检,在调用附件等通道前先判断负载是否可行,从源头避免
工具超时;三是预置常用技能依赖,把 PPT 生成这类高频技能的运行环境提前
装好,避免运行时现场安装失败。这三项都是平台工程动作,不依赖模型升级,
投入产出比最高。

中期做好三件事,把"需要复核"从被动救火变成主动设计:展示知识库解析
覆盖率,让用户一眼看到结论基于多少材料;优先发送链接而非硬塞附件,绕开
附件通道的硬阻塞;将人工确认集中在高风险节点,而不是让用户全程盯守。

长期才谈行业化:固定字段、来源白名单、口径模板与观点版本管理,对应
MuDABench 的四题口径分歧——通用模型之所以在这些题上失分,是因为缺少
行业化的判定规则。到这一步,产品形态就接近真正的垂直行业 Research
Copilot,这也是本项目认为最值得期待的终局形态。

\section{对研究问题的回答}

回到开头的五层框架,每一层都得到了实证回答。

产品能力层面:WorkBuddy 的平台上限不低,换强模型后研究能力接近第一梯队,
但默认配置拖累了实际体验——能力是存在的,问题在于如何被配置。

真实工作流层面:它在材料消化、文件编辑与知识沉淀三个环节创造了价值,在
分发环节制造了摩擦。价值与摩擦并存,而非一边倒,这比"有没有用"的二元
判断更接近真相。

可量化价值层面:检索与知识库场景的 ROI 为正(单次 +240 至 530 元),分发
场景为负(单次 -150 至 -220 元)。省时间与费时间的分界非常清晰,而且
可以量化。

付费意愿层面:个人用户只在检索类场景愿意付费,企业用户价值主张最强但被
稳定性挡在门外。付费意愿的分布与可量化价值的分布高度一致,这本身说明
用户的决策是理性的。

生态放大层面:生态连接是真实的,但价值必须按环节评估——知识库环节放大
了价值,分发环节反而缩小了价值。这是本项目对"生态假设"最重要的修正:
生态不是一个整体变量,而是分环节变量;评估它,必须拆到每个具体工作流
去看。

\section{假设验证汇总}

第 2 章提出的六个假设,到这里可以逐一给出结论。它们不是装饰,而是整项
研究的问题骨架——每一条都有对应的实验证据,或证实、或证伪、或被精确化。

\textbf{H1 工作流假设(部分成立)。}多步骤任务确实能产出交付物:测试一、
测试三与 R1/R2 都完成了多步骤闭环。但 R3 证明多步骤不等于高价值——60
分钟、四次干预、零交付。价值上限成立,可靠性不足,这是 H1 被现实划出的
边界。

\textbf{H2 连续性假设(成立,验证最充分)。}R1 的增量修订、R2 的知识库
版本区分都成功完成,证明 Agent 能承载"持续维护"型工作。研究工作的核心
价值在更新而非首写,H2 是整个研究里被验证最扎实的假设,也直接支撑了
Research Copilot 的构想。

\textbf{H3 用户分层假设(成立)。}个人用户只触达数据层(R2 正 ROI),企微
入口层个人不可达、天然属于企业;MuDABench 的复核需求印证非技术用户需要
复核闭环。用户差异不是营销话术,而是连接深度的客观结果。

\textbf{H4 生态假设(部分成立,需精确化)。}减步骤的环节成立(ima 检索替代
翻 16 份 PDF),读写环节成立(R1/R2);但远程触发个人不可达(企微),分发
环节反而增步骤(R3)。结论:生态价值按环节分化,不能整体断言。

\textbf{H5 商业化假设(成立)。}个人为便利付费(R2 正 ROI),企业为权限、
知识库、安全付费(企微企业专属),付费意愿与实测价值的分布高度一致。

\textbf{H6 专业化假设(强成立)。}MuDABench 四题口径分歧证明:通用模型缺
行业模板就会失分,垂直行业 Research Copilot 需要口径预设与来源规范。
这是整个研究指向最明确的结论,也是优化路线第三步的依据。

六个假设全部得到了实证回应:两个强成立(H2/H6),两个成立(H3/H5),
两个部分成立需精确化(H1/H4)。没有假设被证伪到推翻的程度,但边界条件
被清晰划出——这本身是本研究最有价值的修正。

\section{局限性}

最后需要再次声明本研究的局限,它们共同划定了结论的适用范围。

其一,测试三仅两题,办公编辑的结论基于有限样本,置信度低于其他两项测试;
其二,R1 实际后端为金山文档而非腾讯文档,因此"腾讯文档闭环"未被直接
验证,只能以"文档云闭环"表述,平台差异的结论有待在腾讯文档环境下复测;
其三,L3 的人工耗时为估算区间而非同条件实测,摩擦对比仅用于量级判断,
不构成精确测量;其四,企业微信与企业权限未在真实组织环境下测试,入口层
与协作层的结论仍是空白;其五,WorkBuddy 的换模型复测未覆盖全部任务,
配置敏感性的结论主要建立在测试一上。

后续研究应优先重复知识库与分发两个环节——前者是唯一完整的正价值闭环,
需要更多样本确认其稳定性;后者是最大的负价值环节,需要在工程修复后验证
改善程度。同时应在企业账号环境下补做入口层与协作测试,补齐生态验证的
最后一块拼图。承认这些边界,与坚持前面的结论同样重要:它们不是报告的
瑕疵,而是证据纪律的一部分,让每一个结论都明确标注了自己的适用范围。

% ---------------------------------------------------------------------------
% 参考文献
% ---------------------------------------------------------------------------
\renewcommand{\bibname}{参考文献}
\bibliographystyle{plainnat}
\bibliography{references}

% ---------------------------------------------------------------------------
% 附录
% ---------------------------------------------------------------------------
\ifIncludeAppendix
  \addchap{附录:实验材料与数据说明}

  \section*{A.1 材料存放}
  本项目全部实验材料、任务包冻结说明、评分锚点与原始输出,按模块存放在
  项目仓库的对应目录中;腾讯生态部分按能力层、真实闭环层与摩擦对比层
  三级组织。所有任务包在执行前冻结,版本变更均有记录。

  \section*{A.2 主要数据来源}
  测试一使用公开检索资料,评分含 30 条分层事实抽样、12 个核心模型家族
  检查表与每家 20 条引用抽查;测试二使用 Factcheck-GPT、FactBench 与
  FACTS Grounding 公开样本;测试三使用 SpreadsheetBench 2 样本;腾讯生态
  测试使用 HERB 材料包与本项目自建的 MuDABench 中文企业资料库。各基准的
  原始引用见参考文献。

  \section*{A.3 评分标准速查}
  测试一按五个维度加权(事实可靠性、研究完整性、分析观点、来源与引用、
  效率);测试二按两个子任务分别评分后取平均;测试三直接核验交付文件
  (可打开性、公式正确性、图表对象真实性);腾讯生态部分按各任务的通过
  判据核验,通过率等于通过项数除以适用项数。

  \section*{A.4 假设与证据对应}
  六大假设(H1 工作流、H2 连续性、H3 用户分层、H4 生态、H5 商业化、
  H6 专业化)的完整定义见正文第 2 章"研究假设"一节,逐条验证结论见
  正文第 14 章"假设验证汇总"一节。附录不再重复,以免与正文脱节。
\fi

% 页脚总页数标记
\phantomsection
\label{LastNumberedPage}
\clearpage
\pagestyle{empty}

\end{document}
